This study has a few limitations. First, we only evaluate on AG News, so our conclusions may not transfer to datasets with different writing styles or label structures. A second limitation is our choice of a 5{,}000-item BPE vocabulary. This is a reasonable size for a small study, but we did not check token coverage on the test set, so we cannot say how much information may have been lost because of vocabulary size. A larger vocabulary might have preserved more useful distinctions, especially for named entities and domain-specific terms.

Another important limitation is the tuning setup. Hyperparameter search was done with only 3 epochs per configuration, and we did not explicitly verify that this was long enough for strong and weak configurations to separate clearly. This means some settings may have looked worse simply because they learned more slowly, not because they were truly worse. The CNN-128 result also shows that small differences between runs should be interpreted carefully.

Regularisation is another place where our choices may have been too limited. In the best configurations, weight decay was either very small or zero, and the learning curves suggest that the LSTMs overfit slightly more than the CNNs. Because of this, it would have been useful to test higher weight decay values more aggressively, especially for the LSTM models.

Finally, we compare only CNN and LSTM models trained from scratch, so this is a comparison between two classical neural architectures rather than a comparison with the current strongest text-classification models. That does not affect the internal comparison in this report, but it does limit how far the absolute performance numbers can be taken.
